\scanpage{8}
Notons que pour une forme bilinéaire
\[
 \varphi(x,y)=\sum_{1\leq i,j\leq3}\varphi_{ij}x_i y_j
\]
on a le discriminant de $\varphi$ (ou de la forme quadratique $q$, si $\varphi$ est associé à une telle forme)
\[
 \delta(\varphi)=\unsure{-|}\det(\varphi_{ij})\,\omega'^{\unsure{\otimes}2}
 \qquad(\text{où }\omega'=e'_1\wedge e'_{\unsure{1|2}}\wedge e'_3), \tag*{(27)}
\]
et si $\varphi$ est associé à une forme quadratique
\[
 \inserted{q(x,y,z)=}ax^2+by^2+cz^2+pyz+qzx+rxy \tag*{(28)}
\]
de sorte que la matrice de $\varphi$ est
\[
 \varphi=\begin{pmatrix}2a&r&q\\r&2b&p\\q&p&2c\end{pmatrix} \tag*{(29)}
\]
on a
\[
 \delta(q)=\delta(\varphi)=2(4abc+pqr-ap^2-bq^2-cr^2)\,\omega'^{\otimes2} \tag*{(30)}
\]
et le ``discriminant divisé'' $\delta'(q)$ de $q$ est donné par
\[
 \delta'(q)=(4abc+pqr-ap^2-bq^2-cr^2)\,\omega'^{\otimes2}, \tag*{(31)}
\]
\deleted{ce qui} Dans le cas de la forme $q$ de (21) d'où (23 bis), on trouve donc
\[
 \delta'(q)=-\omega'^{\otimes2}\qquad(\omega'=e'_1\wedge e'_2\wedge e'_3) \tag*{(32)}
\]

En s'inspirant de cette circonstance, on définit une structure \uline{spéciale orthogonale} sur un Module \deleted{\unsure{quadrati-|quadratique}} vectoriel de rang 3, comme la donnée d'un couple $(q,\omega)$, où $q$ est une forme quadratique sur $E$, $\omega$ une base de $\underline{\det} E$, satisfaisant

% archive sheet number: 7
